\chapter{2014年新课标I卷（理）}

\section{单选题}

\begin{problem}
已知集合 \(A = \{x \mid x^2 - 2x - 3 \geqslant 0\}\)，\(B = \{x \mid -2 \leqslant x < 2\}\)，则 \(A \cap B =\)（\quad）

\choices
    {\([-2, -1]\)}
    {\([-1, 2)\)}
    {\([-1, 1]\)}
    {\([1,2)\)}
\end{problem}

\begin{answer}
A
\end{answer}

\begin{solution}
由 \(x^2-2x-3=(x-3)(x+1)\geqslant0\)，得
\(A=(-\infty,-1]\cup[3,+\infty)\). 因为 \(B=[-2,2)\)，所以
\(A\cap B=[-2,-1]\)，故选 A.
\end{solution}

\begin{problem}
\(\frac{(1 + \i)^3}{(1 - \i)^2} =\)（\quad）

\choices
    {\(1 + \i\)}
    {\(1 - \i\)}
    {\(-1 + \i\)}
    {\(-1 - \i\)}
\end{problem}

\begin{answer}
D
\end{answer}

\begin{solution}
\[
(1+\i)^3=(1+\i)^2(1+\i)=2\i(1+\i)=-2+2\i
\]
\[
(1-\i)^2=-2\i
\]
因此
\[
\frac{(1+\i)^3}{(1-\i)^2}=\frac{-2+2\i}{-2\i}=-1-\i
\]
故选 D.
\end{solution}

\begin{problem}
设函数 \(f(x)\)，\(g(x)\) 的定义域都为 \( \R \)，且 \( f(x) \) 是奇函数，\( g(x) \) 是偶函数，则下列结论正确的是（\quad）

\choices
    {\(f(x)g(x)\) 是偶函数}
    {\(|f(x)|g(x)\) 是奇函数}
    {\(|g(x)|f(x)\) 是奇函数}
    {\(|f(x)g(x)|\) 是奇函数}
\end{problem}

\begin{answer}
C
\end{answer}

\begin{solution}
由 \(f(-x)=-f(x)\)，\(g(-x)=g(x)\)，可得
\(|f(-x)|=|f(x)|\)，\(|g(-x)|=|g(x)|\). 因而
\[
|g(-x)|f(-x)=|g(x)|[-f(x)]=-\lvert g(x)\rvert f(x)
\]
所以 \(|g(x)|f(x)\) 是奇函数，选 C. 另外，\(f(x)g(x)\) 是奇函数，
\(|f(x)|g(x)\) 和 \(|f(x)g(x)|\) 是偶函数，故其余选项均不正确.
\end{solution}

\begin{problem}
已知 \(F\) 是双曲线 \(C: x^{2} - my^{2} = 3m \) （\(m > 0\)）的一个焦点，则点 \(F\) 到 \(C\) 的一条渐近线的距离为（\quad）

\choices
    {\(\sqrt{3}\)}
    {3}
    {\(\sqrt{3} m\)}
    {\(3m\)}
\end{problem}

\begin{answer}
A
\end{answer}

\begin{solution}
将双曲线方程化为
\[
\frac{x^2}{3m}-\frac{y^2}{3}=1
\]
所以 \(a^2=3m\)，\(b^2=3\)，焦距参数满足
\(c^2=a^2+b^2=3(m+1)\). 一条渐近线可写为
\(x-\sqrt m y=0\). 因为 \(F\) 的坐标为 \((\pm c,0)\)，点 \(F\) 到该直线的距离为
\[
\frac{|\pm c|}{\sqrt{1+m}}=\frac{\sqrt{3(m+1)}}{\sqrt{m+1}}=\sqrt3
\]
故选 A.
\end{solution}

\begin{problem}
\(4\) 位同学各自在周六，周日两天中任选一天参加公益活动，则周六，周日都有同学参加公益活动的概率为（\quad）

\choices
    {\( \frac{1}{8} \)}
    {\( \frac{3}{8} \)}
    {\( \frac{5}{8} \)}
    {\( \frac{7}{8} \)}
\end{problem}

\begin{answer}
D
\end{answer}

\begin{solution}
每位同学有两种等可能选择，因此所有选择共有 \(2^4=16\) 种. 周六、周日都有同学参加的对立事件是四人全选周六或四人全选周日，
共有 \(2\) 种. 所求概率为
\[
1-\frac{2}{16}=\frac78
\]
故选 D.
\end{solution}

\begin{problem}
如图，圆 \(O\) 的半径为 \(1\)，\(A\) 是圆上的定点，\(P\) 是圆上的动点，角 \(x\) 的始边为射线 \(OA\)，终边为射线 \(OP\)，过点 \(P\) 作直线 \(OA\) 的垂线，垂足为 \(M\)，将点 \(M\) 到直线 \(OP\) 的距离表示为 \(x\) 的函数 \(f(x)\)，则 \(y = f(x)\) 在 \(\left[0, \pi\right]\) 上的图象大致为（\quad）

\bitmapfigure[max width=0.62\linewidth]{img/2014/new_curriculum_paper_1_science/q6_fig2.png}

\choices
    {\choicebitmap[width=0.15\paperwidth]{img/2014/new_curriculum_paper_1_science/q6_optA.png}}
    {\choicebitmap[width=0.15\paperwidth]{img/2014/new_curriculum_paper_1_science/q6_optB.png}}
    {\choicebitmap[width=0.15\paperwidth]{img/2014/new_curriculum_paper_1_science/q6_optC.png}}
    {\choicebitmap[width=0.15\paperwidth]{img/2014/new_curriculum_paper_1_science/q6_optD.png}}
\end{problem}

\begin{answer}
C
\end{answer}

\begin{solution}
建立直角坐标系，使 \(O=(0,0)\)，\(A=(1,0)\). 当 \(x\in[0,\pi]\) 时，
\(P=(\cos x,\sin x)\)，而垂足为 \(M=(\cos x,0)\). 直线 \(OP\) 的方向向量为
\((\cos x,\sin x)\)，且其长度为 \(1\)，所以点 \(M\) 到直线 \(OP\) 的距离为
\[
f(x)=|\cos x\sin x|
\]
在 \([0,\pi]\) 上，函数在 \(0\)，\(\dfrac\pi2\)，\(\pi\) 处为零，在
\(\left(0,\dfrac\pi2\right)\) 和 \(\left(\dfrac\pi2,\pi\right)\) 内均为正，且图象关于
\(x=\dfrac\pi2\) 对称. 因此图象为 C.
\end{solution}

\begin{problem}
执行如图的程序框图，若输入的 \(a\)，\(b\)，\(k\) 分别为 \(1\)，\(2\)，\(3\)，则输出的 \(M =\)（\quad）

\bitmapfigure[max width=0.56\linewidth]{img/2014/new_curriculum_paper_1_science/q7_fig1.png}

\choices
    {\(\frac{20}{3}\)}
    {\(\frac{7}{2}\)}
    {\(\frac{16}{5}\)}
    {\(\frac{15}{8}\)}
\end{problem}

\begin{answer}
D
\end{answer}

\begin{solution}
初始时 \(n=1\)，且
\(M=a+\dfrac1b=1+\dfrac12=\dfrac32\). 随后依次执行
\(a=b\)、\(b=M\)，得到 \(a=2\)、\(b=\dfrac32\)，再令 \(n=2\).

当 \(n=2\leqslant3\) 时，
\(M=2+\dfrac{1}{3/2}=\dfrac83\). 随后执行
\(a=b\)、\(b=M\)，得到 \(a=\dfrac32\)、\(b=\dfrac83\)，再令 \(n=3\).

当 \(n=3\leqslant3\) 时，
\(M=\dfrac32+\dfrac{1}{8/3}=\dfrac{15}{8}\). 随后执行
\(a=b\)、\(b=M\)，再令 \(n=4\). 此时 \(n\leqslant k\) 不成立，程序输出
\(M=\dfrac{15}{8}\)，故选 D.
\end{solution}

\begin{problem}
设 \(\alpha \in \left(0, \frac{\pi}{2}\right)\)，\(\beta \in \left(0, \frac{\pi}{2}\right)\)，且 \(\tan \alpha = \frac{1 + \sin \beta}{\cos \beta}\)，则（\quad）

\choices
    {\(3\alpha -\beta = \frac{\pi}{2}\)}
    {\(3\alpha +\beta = \frac{\pi}{2}\)}
    {\(2\alpha -\beta = \frac{\pi}{2}\)}
    {\(2\alpha +\beta = \frac{\pi}{2}\)}
\end{problem}

\begin{answer}
C
\end{answer}

\begin{solution}
利用半角公式，
\[
\frac{1+\sin\beta}{\cos\beta}
 =\frac{\left(\sin\dfrac\beta2+\cos\dfrac\beta2\right)^2}
 {\left(\cos\dfrac\beta2-\sin\dfrac\beta2\right)
  \left(\cos\dfrac\beta2+\sin\dfrac\beta2\right)}
 =\frac{\cos\dfrac\beta2+\sin\dfrac\beta2}
 {\cos\dfrac\beta2-\sin\dfrac\beta2}
 =\tan\left(\frac\pi4+\frac\beta2\right)
\]
由于 \(\alpha\)，\(\beta\) 均在 \(\left(0,\dfrac\pi2\right)\) 内，等式两边对应的角均在
\(\left(0,\dfrac\pi2\right)\) 内，故
\(\alpha=\dfrac\pi4+\dfrac\beta2\). 从而
\(2\alpha-\beta=\dfrac\pi2\)，选 C.
\end{solution}

\begin{problem}
不等式组 \(\begin{cases}
x + y \geqslant 1,\\
x - 2y \leqslant 4
\end{cases}\) 的解集记为 \(D\). 有下面四个命题：

\begin{circlelist}
\item \(p_1\colon \forall (x,y)\in D\)，\(x + 2y\geqslant -2\)；
\item \(p_2\colon \exists (x,y)\in D\)，\(x + 2y\geqslant 2\)；
\item \(p_3\colon \forall (x,y)\in D\)，\(x + 2y\leqslant 3\)；
\item \(p_4\colon \exists (x,y)\in D\)，\(x + 2y\leqslant -1\).
\end{circlelist}

其中真命题是（\quad）

\choices
    {\(p_2, p_3\)}
    {\(p_1\)，\(p_2\)}
    {\(p_1\)，\(p_4\)}
    {\(p_1\)，\(p_3\)}
\end{problem}

\begin{answer}
B
\end{answer}

\begin{solution}
由 \(x+y\geqslant1\) 和 \(x-2y\leqslant4\)，可得
\(x\geqslant1-y\)，\(x\leqslant4+2y\)，所以必须有 \(y\geqslant-1\). 因而
\[
x+2y\geqslant(1-y)+2y=1+y\geqslant0
\]
所以对所有 \((x,y)\in D\)，均有 \(x+2y\geqslant-2\)，命题 \(p_1\) 为真，
而 \(p_4\) 为假.

取 \((x,y)=(3,0)\in D\)，有 \(x+2y=3\geqslant2\)，故 \(p_2\) 为真. 另一方面，
取 \((x,y)=(4,0)\in D\)，有 \(x+2y=4>3\)，故 \(p_3\) 为假. 真命题为
\(p_1\)，\(p_2\)，选 B.
\end{solution}

\begin{problem}
已知抛物线 \(C: y^{2} = 8x\) 的焦点为 \(F\)，准线为 \(l\)，\(P\) 是 \(l\) 上一点，\(Q\) 是直线 \(PF\) 与 \(C\) 的一个交点，若 \(\overrightarrow{FP} = 4\overrightarrow{FQ}\)，则 \(|QF| =\)（\quad）

\choices
    {\(\frac{7}{2}\)}
    {3}
    {\(\frac{5}{2}\)}
    {2}
\end{problem}

\begin{answer}
B
\end{answer}

\begin{solution}
抛物线可写为 \(y^2=4px\)，故 \(p=2\)，焦点为 \(F=(2,0)\)，准线为
\(x=-2\). 设 \(P=(-2,t)\). 由
\(\overrightarrow{FP}=4\overrightarrow{FQ}\)，得
\[
Q=F+\frac14(P-F)=\left(1,\frac t4\right)
\]
因为 \(Q\) 在抛物线上，
\[
\left(\frac t4\right)^2=8\times1
\]
从而 \(Q=(1,\pm2\sqrt2)\). 因此
\[
|QF|=\sqrt{(1-2)^2+(\pm2\sqrt2)^2}=3
\]
故选 B.
\end{solution}

\begin{problem}
已知函数 \( f(x) = ax^3 - 3x^2 + 1 \)，若 \( f(x) \) 存在唯一的零点 \( x_0 \)，且 \( x_0 > 0 \)，则 \( a \) 的取值范围为（\quad）

\choices
    {\((2, +\infty)\)}
    {\((1, +\infty)\)}
    {\((- \infty, -2)\)}
    {\((- \infty, -1)\)}
\end{problem}

\begin{answer}
C
\end{answer}

\begin{solution}
先考虑 \(a=0\). 此时 \(f(x)=1-3x^2\)，有两个零点
\(x=\pm\dfrac{\sqrt3}{3}\)，不符合题意.

再考虑 \(a>0\). 由于
\(\displaystyle\lim_{x\to-\infty}f(x)=-\infty\)，而 \(f(0)=1>0\)，由零点存在定理，在负半轴上至少有一个零点，不可能满足“唯一的零点为正数”.

最后考虑 \(a<0\). 此时
\[
f'(x)=3x(ax-2)
\]
函数在 \(x=\dfrac2a\) 和 \(x=0\) 处取得极值，且
\(f(0)=1\)，\(f\left(\dfrac2a\right)=1-\dfrac4{a^2}\).
当 \(x<0\) 时，\(f'(x)\) 在 \(x=\dfrac2a\) 左侧为负、在
\(\left(\dfrac2a,0\right)\) 内为正，因此 \(x=\dfrac2a\) 是负半轴上的极小值点.
又因为 \(f(x)\to+\infty\)（\(x\to-\infty\)），所以负半轴没有零点的充要条件是该极小值大于 \(0\)，即
\[
1-\frac4{a^2}>0
\]
由 \(a<0\)，得 \(a<-2\). 对 \(x>0\)，有 \(f'(x)<0\)，且
\(f(0)=1>0\)、\(f(x)\to-\infty\)（\(x\to+\infty\)），所以正半轴恰有一个零点.
因此唯一零点为正数的充要条件是 \(a<-2\)，选 C.
\end{solution}

\begin{problem}
如图，网格纸上小正方形的边长为 \(1\)，粗实线画出的是某多面体的三视图，则该多面体的各条棱中，最长的棱的长度为（\quad）

\bitmapfigure[max width=0.42\linewidth]{img/2014/new_curriculum_paper_1_science/q12_fig1.png}

\choices
    {\(6\sqrt{2}\)}
    {6}
    {\(4\sqrt{2}\)}
    {4}
\end{problem}

\begin{answer}
B
\end{answer}

\begin{solution}
由三个视图可还原出一个三棱锥 \(D-ABC\)：底面 \(ABC\) 为等腰直角三角形，网格给出
\(AB\perp BC\) 且 \(AB=BC=4\)；侧视图给出点 \(D\) 在底面上的射影 \(M\) 是
\(BC\) 的中点，且 \(DM\perp\) 底面、\(DM=4\). 取底面为 \(z=0\)，建立直角坐标系
\(A=(0,0,0)\)，\(B=(4,0,0)\)，\(C=(4,4,0)\)，于是
\(M=(4,2,0)\)，\(D=(4,2,4)\).
于是六条棱的长度分别为
\(AB=BC=4\)，\(AC=\sqrt{4^2+4^2}=4\sqrt2\)，
\(BD=CD=\sqrt{2^2+4^2}=2\sqrt5\)，\(AD=\sqrt{4^2+2^2+4^2}=6\).
比较它们的平方，\(36>32>20>16\)，所以最长棱长为 \(6\)，选 B.
\end{solution}

\section{填空题}

\begin{problem}
\((x - y)(x + y)^8\) 的展开式中 \(x^{2}y^{7}\) 的系数为\fillinblank{}.
\end{problem}

\begin{answer}
\(-20\)
\end{answer}

\begin{solution}
展开式中，取前面的 \(x\) 与 \((x+y)^8\) 中的
\(\mathrm C_8^1xy^7\)，得到系数 \(\mathrm C_8^1=8\). 取前面的 \(-y\) 与
\((x+y)^8\) 中的 \(\mathrm C_8^2x^2y^6\)，得到系数
\(-\mathrm C_8^2=-28\). 因此 \(x^2y^7\) 的系数为
\[
8-28=-20
\]
\end{solution}

\begin{problem}
甲，乙，丙三位同学被问到是否去过 \(A\)，\(B\)，\(C\) 三个城市时，甲说：我去过的城市比乙多，但没去过 \(B\) 城市；乙说：我没去过 \(C\) 城市；丙说：我们三人去过同一个城市. 由此可判断乙去过的城市为\fillinblank{}.
\end{problem}

\begin{answer}
\(A\)
\end{answer}

\begin{solution}
甲没有去过 \(B\)，乙没有去过 \(C\)，而三人去过同一个城市，所以这个共同去过的城市只能是 \(A\). 因此乙一定去过 \(A\).

若乙还去过 \(B\)，则乙去过两个城市. 甲没有去过 \(B\)，在三个城市中至多去过 \(A\)、\(C\) 两个城市，不可能比乙去过的城市多，矛盾. 因而乙没有去过 \(B\)，结合乙没有去过 \(C\)，可知乙去过的城市恰为 \(A\).
\end{solution}

\begin{problem}
已知 \(A\)，\(B\)，\(C\) 是圆 \(O\) 上的三点，若 \( \overrightarrow{AO} = \frac{1}{2}\left(\overrightarrow{AB} + \overrightarrow{AC}\right) \)，则 \( \overrightarrow{AB} \) 与 \( \overrightarrow{AC} \) 的夹角为\fillinblank{}.
\end{problem}

\begin{answer}
\(90^{\circ}\)
\end{answer}

\begin{solution}
由向量中点公式，
\[
\overrightarrow{AO}=\frac12\left(\overrightarrow{AB}+\overrightarrow{AC}\right)
\]
说明点 \(O\) 是线段 \(BC\) 的中点. 又因为 \(B\)，\(C\) 在以 \(O\) 为圆心的圆上，
\(OB=OC\)，所以 \(BC\) 是圆 \(O\) 的直径. 因此由圆周角定理，
\(\angle BAC=90^{\circ}\)，即 \(\overrightarrow{AB}\) 与 \(\overrightarrow{AC}\) 的夹角为
\(90^{\circ}\).
\end{solution}

\begin{problem}
已知 \(a\)，\(b\)，\(c\) 分别为 \(\triangle ABC\) 的三个内角 \(A\)，\(B\)，\(C\) 的对边. \(a = 2\)，且 \((2 + b)(\sin A - \sin B) = (c - b)\sin C\)，则 \(\triangle ABC\) 面积的最大值为\fillinblank{}.
\end{problem}

\begin{answer}
\(\sqrt{3}\)
\end{answer}

\begin{solution}
由正弦定理，存在外接圆半径 \(R\)，使
\(\sin A=\dfrac{a}{2R}\)，\(\sin B=\dfrac{b}{2R}\)，\(\sin C=\dfrac{c}{2R}\). 代入题设得
\[
\frac{(2+b)(2-b)}{2R}=\frac{c(c-b)}{2R}
\]
即
\[
b^2+c^2-bc=4
\]
另一方面，余弦定理给出
\(4=b^2+c^2-2bc\cos A\). 两式相减，得
\(bc(2\cos A-1)=0\). 因为 \(b\)，\(c>0\)，所以 \(\cos A=\dfrac12\)，即
\(A=60^{\circ}\).

三角形面积为
\[
S=\frac12bc\sin A=\frac{\sqrt3}{4}bc
\]
又
\[
4=b^2+c^2-bc=(b-c)^2+bc\geqslant bc
\]
所以 \(bc\leqslant4\)，从而 \(S\leqslant\sqrt3\). 当 \(b=c=2\) 时等号成立，
故面积的最大值为 \(\sqrt3\).
\end{solution}

\section{解答题}

\begin{problem}
已知数列 \(\{a_{n}\}\) 的前 \(n\) 项和为 \(S_{n}\)，\(a_{1} = 1\)，\(a_{n} \neq 0\)，\(a_{n}a_{n+1} = \lambda S_{n} - 1\)，其中 \(\lambda\) 为常数.

\begin{enumerate}
\item 证明：\(a_{n+2}-a_{n}=\lambda\)；
\item 是否存在 \(\lambda\)，使得 \(\{a_{n}\}\) 为等差数列？并说明理由.
\end{enumerate}
\end{problem}

\begin{solution}
\begin{enumerate}
\item 由题设，对任意正整数 \(n\)，有
\[
a_na_{n+1}=\lambda S_n-1
\]
将 \(n\) 换成 \(n+1\)，并利用 \(S_{n+1}=S_n+a_{n+1}\)，得
\[
a_{n+1}a_{n+2}=\lambda S_{n+1}-1
 =\lambda S_n+\lambda a_{n+1}-1
\]
两式相减，得到
\[
a_{n+1}(a_{n+2}-a_n)=\lambda a_{n+1}
\]
由于 \(a_{n+1}\ne0\)，所以
\(a_{n+2}-a_n=\lambda\).

\item 取 \(n=1\) 代入已知递推式，得
\(a_1a_2=a_2=\lambda S_1-1=\lambda-1\)，即 \(a_2=\lambda-1\).
若数列为等差数列，则其公差为
\(d=a_2-a_1=\lambda-2\)，从第（1）问又有
\(a_3-a_1=\lambda\)，所以
\[
a_3=1+\lambda
\]
另一方面，等差数列应满足
\(a_3=a_1+2d=1+2(\lambda-2)=2\lambda-3\). 比较两式得
\(\lambda=4\).

当 \(\lambda=4\) 时，由 \(a_1=1\)、\(a_2=3\) 及
\(a_{n+2}-a_n=4\)，可得 \(a_n=2n-1\). 此时
\(S_n=1+3+\cdots+(2n-1)=n^2\)，且
\[
a_na_{n+1}=(2n-1)(2n+1)=4n^2-1=4S_n-1
\]
满足题设条件，故存在且唯一的 \(\lambda=4\)，数列为等差数列.
\end{enumerate}
\end{solution}

\begin{problem}
从某企业生产的某种产品中抽取 \(500\) 件，测量这些产品的一项质量指标值，由测量结果得如下频率分布直方图：

\begin{enumerate}
\item 求这 \(500\) 件产品质量指标值的样本平均数 \( \bar{x} \) 和样本方差 \( s^{2} \) （同一组数据用该区间的中点值作代表）；
\item 由频率分布直方图可以认为，这种产品的质量指标值 \(Z\) 服从正态分布 \(N(\mu, \sigma^2)\)，其中 \(\mu\) 近似为样本平均数 \(\bar{x}\)，\(\sigma^2\) 近似为样本方差 \(s^2\).

\begin{enumerate}
\item 利用该正态分布，求 \( P(187.8 < Z < 212.2) \)；
\item 某用户从该企业购买了 \(100\) 件这种产品，记 \(X\) 表示这 \(100\) 件产品中质量指标值位于区间 \((187.8, 212.2)\) 的产品件数，利用（i）的结果，求 \(E(X)\).
\end{enumerate}
\end{enumerate}

附：\( \sqrt{150} \approx 12.2 \)，若 \( Z \sim N(\mu, \sigma^{2}) \)，则 \( P(\mu - \sigma < Z < \mu + \sigma) = 0.6826 \)，\( P(\mu - 2\sigma < Z < \mu + 2\sigma) = 0.9544 \).

\begin{center}
\texfigure[float=inline,max width=0.60\linewidth]{img_repaint/2014/new_curriculum_paper_1_science/q18_fig1.tex}
\end{center}
\FloatBarrier
\end{problem}

\begin{solution}
\begin{enumerate}
\item 由直方图可读出区间 \([165,175)\)，\([175,185)\)，\([185,195)\)，
\([195,205)\)，\([205,215)\)，\([215,225)\)，\([225,235)\) 的组中值分别为
\(170\)，\(180\)，\(190\)，\(200\)，\(210\)，\(220\)，\(230\)，频率分别为
\(0.02\)，\(0.09\)，\(0.22\)，\(0.33\)，\(0.24\)，\(0.08\)，\(0.02\). 这些频率之和为 \(1\). 因此样本平均数为
\[
\bar{x}=170\times0.02+180\times0.09+190\times0.22+200\times0.33
 +210\times0.24+220\times0.08+230\times0.02=200
\]
以组中值代表该组数据，样本方差为
\[
\begin{aligned}
s^2={}&(170-200)^2\times0.02+(180-200)^2\times0.09
 +(190-200)^2\times0.22\\
&+(200-200)^2\times0.33+(210-200)^2\times0.24
 +(220-200)^2\times0.08+(230-200)^2\times0.02\\
={}&150.
\end{aligned}
\]

\item
\begin{enumerate}
\item 由 \(\mu\approx\bar{x}=200\)，\(\sigma^2\approx s^2=150\)，得
\(\sigma\approx\sqrt{150}\approx12.2\). 因而
\(187.8=200-12.2\approx\mu-\sigma\)，\(212.2=200+12.2\approx\mu+\sigma\).
所以
\[
P(187.8<Z<212.2)=P(\mu-\sigma<Z<\mu+\sigma)=0.6826
\]
\item 每件产品的质量指标值落在区间 \((187.8,212.2)\) 内的概率为
\(p=0.6826\). 因而
\[
E(X)=100p=100\times0.6826=68.26
\]
\end{enumerate}
\end{enumerate}
\end{solution}

\begin{problem}
如图，三棱柱 \(ABC - A_{1}B_{1}C_{1}\) 中，侧面 \(BB_{1}C_{1}C\) 为菱形，\(AB \perp B_{1}C\).

\begin{enumerate}
\item 证明：\(AC = AB_{1}\)；
\item 若 \(AC \perp AB_{1}\)，\(\angle CBB_{1} = 60^{\circ}\)，\(AB = BC\)，求二面角 \(A - A_{1}B_{1} - C_{1}\) 的余弦值.
\end{enumerate}

\bitmapfigure[max width=0.62\linewidth]{img/2014/new_curriculum_paper_1_science/q19_fig1.png}
\end{problem}

\begin{solution}
\begin{enumerate}
\item 设 \(\overrightarrow{AB}=\bs w\)，\(\overrightarrow{BC}=\bs u\)，
\(\overrightarrow{BB_1}=\bs v\). 因为四边形 \(BB_1C_1C\) 为菱形，所以
\(|\bs u|=|\bs v|\). 又因为棱柱的侧棱平行且相等，
\(\overrightarrow{AB_1}=\bs w+\bs v\)，\(\overrightarrow{AC}=\bs w+\bs u\).
由 \(AB\perp B_1C\) 得
\(\bs w\cdot(\bs u-\bs v)=0\). 因此
\[
\begin{aligned}
AC^2-AB_1^2
&=|\bs w+\bs u|^2-|\bs w+\bs v|^2\\
&=2\bs w\cdot(\bs u-\bs v)+|\bs u|^2-|\bs v|^2=0.
\end{aligned}
\]
所以 \(AC=AB_1\).

\item 取 \(BC=BB_1=1\)，建立空间直角坐标系，令
\(B=(0,0,0)\)，\(C=(1,0,0)\)，\(B_1=\left(\dfrac12,\dfrac{\sqrt3}{2},0\right)\).
设 \(\overrightarrow{AB}=\bs w=(x,y,z)\). 由 \(AB=BC=1\)、
\(AB\perp B_1C\) 以及 \(AC\perp AB_1\)，分别得到
\(x^2+y^2+z^2=1\)，\(\bs w\cdot(C-B_1)=0\)，
\((\bs w+C-B)\cdot(\bs w+B_1-B)=0\).
由于 \(C-B_1=\left(\frac12,-\frac{\sqrt3}{2},0\right)\)，由
\(\bs w\cdot(C-B_1)=0\) 得 \(x=\sqrt3y\). 又因为
\((C-B)\cdot(B_1-B)=\frac12\)，将 \(x=\sqrt3y\) 代入
\((\bs w+C-B)\cdot(\bs w+B_1-B)=0\)，得
\[
1+2\sqrt3y+\frac12=0
\]
所以 \(y=-\dfrac{\sqrt3}{4}\)，\(x=-\dfrac34\). 再由
\(x^2+y^2+z^2=1\) 得 \(z^2=\dfrac14\). 取 \(z\) 轴方向使
\(z=\dfrac12\)，于是
\(\bs w=\left(-\dfrac34,-\dfrac{\sqrt3}{4},\dfrac12\right)\).
于是
\(A=\left(\dfrac34,\dfrac{\sqrt3}{4},-\dfrac12\right)\)，
\(A_1=A+B_1-B=\left(\dfrac54,\dfrac{3\sqrt3}{4},-\dfrac12\right)\)，
且 \(C_1=C+B_1-B=\left(\frac32,\frac{\sqrt3}{2},0\right)\).

平面 \(AA_1B_1\) 的一个法向量和平面 \(C_1A_1B_1\) 的一个法向量可分别取
\(\bs n_1=(A_1-A)\times(B_1-A)\)，
\(\bs n_2=(C_1-B_1)\times(A_1-B_1)\).
计算可得
\(\bs n_1=\left(\dfrac{\sqrt3}{4},-\dfrac14,\dfrac{\sqrt3}{4}\right)\)，
\(\bs n_2=\left(0,\dfrac12,\dfrac{\sqrt3}{4}\right)\).
两法向量的长度均为 \(\dfrac{\sqrt7}{4}\)，且
\(\bs n_1\cdot\bs n_2=\dfrac1{16}\). 因此二面角
\(A-A_1B_1-C_1\) 的余弦值为
\[
\frac{|\bs n_1\cdot\bs n_2|}{|\bs n_1||\bs n_2|}
 =\frac{1/16}{(\sqrt7/4)^2}=\frac17
\]
\end{enumerate}
\end{solution}

\begin{problem}
已知点 \(A(0, -2)\)，椭圆 \(E: \frac{x^2}{a^2} + \frac{y^2}{b^2} = 1 \) （\(a > b > 0\)）的离心率为 \(\frac{\sqrt{3}}{2}\)，\(F\) 是椭圆的右焦点，直线 \(AF\) 的斜率为 \(\frac{2\sqrt{3}}{3}\)，\(O\) 为坐标原点.

\begin{enumerate}
\item 求 \(E\) 的方程；
\item 设过点 \(A\) 的动直线 \(l\) 与 \(E\) 相交于 \(P\)，\(Q\) 两点. 当 \(\triangle OPQ\) 的面积最大时，求 \(l\) 的方程.
\end{enumerate}
\end{problem}

\begin{solution}
\begin{enumerate}
\item 设椭圆的半焦距为 \(c\)，则右焦点为 \(F=(c,0)\). 由直线 \(AF\) 的斜率为
\(\dfrac{2\sqrt3}{3}\)，得
\(\frac{2}{c}=\frac{2\sqrt3}{3}\)，\(qquad c=\sqrt3\).
又由离心率 \(e=\dfrac ca=\dfrac{\sqrt3}{2}\)，得 \(a=2\). 因为
\(b^2=a^2-c^2=1\)，所以椭圆方程为
\[
\frac{x^2}{4}+y^2=1
\]

\item 若 \(l\) 为竖直直线，则 \(l:x=0\)，它与椭圆的交点为
\((0,1)\)、\((0,-1)\)，点 \(O\)、\(P\)、\(Q\) 共线，三角形 \(OPQ\) 的面积为 \(0\)，不可能取得最大值.
因此以下设动直线 \(l\) 的斜率为 \(k\)，则 \(l:y=kx-2\). 令直线上的点为
\(A+t(1,k)=(t,kt-2)\)，与椭圆联立得
\[
\frac{t^2}{4}+(kt-2)^2=1
\]
即
\[
(1+4k^2)t^2-16kt+12=0
\]
设两个交点对应的参数为 \(t_1\)，\(t_2\). 因为
\(\det(A+t_1(1,k),A+t_2(1,k))=2(t_2-t_1)\)，所以
\[
S_{\triangle OPQ}=|t_2-t_1|
\]
由一元二次方程根的判别式，
\[
|t_2-t_1|=\frac{\sqrt{(-16k)^2-4(1+4k^2)\cdot12}}{1+4k^2}
 =\frac{4\sqrt{4k^2-3}}{1+4k^2}
\]
令 \(u=4k^2-3\geqslant0\)，则
\[
S_{\triangle OPQ}=\frac{4\sqrt u}{u+4}\leqslant1
\]
其中利用了 \(u+4\geqslant4\sqrt u\). 等号成立当且仅当 \(u=4\)，即
\(4k^2=7\)，所以
\[
l:y=\frac{\sqrt7}{2}x-2
\quad\text{或}\quad
l:y=-\frac{\sqrt7}{2}x-2
\]
\end{enumerate}
\end{solution}

\begin{problem}
设函数 \( f(x) = a \e^{x} \ln x + \frac{b \e^{x-1}}{x} \)，曲线 \( y = f(x) \) 在点 \( (1, f(1)) \) 处的切线方程为 \( y = \e(x - 1) + 2 \).

\begin{enumerate}
\item 求 \(a\)，\(b\)；
\item 证明：\(f(x) > 1\).
\end{enumerate}
\end{problem}

\begin{solution}
\begin{enumerate}
\item 函数定义域为 \(x>0\). 由切线方程可知
\(f(1)=2\)，且 \(f'(1)=\e\). 直接计算得
\[
f(1)=b
\]
以及
\[
f'(x)=a\e^x(\ln x+1)+b\e^{x-1}\left(\frac1x-\frac1{x^2}\right)
\]
代入 \(x=1\) 得 \(f'(1)=a\e\).
因此 \(b=2\)，\(a=1\).

\item 代入 \(a=1\)，\(b=2\)，得
\[
f(x)=\e^x\ln x+\frac{2\e^{x-1}}x
 =\frac{\e^{x-1}}x\left(\e x\ln x+2\right)
\]
令 \(h(x)=x\ln x\)，则
\[
h'(x)=\ln x+1
\]
函数 \(h(x)\) 在 \(x=\e^{-1}\) 处取得最小值
\(h(\e^{-1})=-\e^{-1}\)，所以
\(\e x\ln x+2\geqslant1\). 另一方面，\(\ln x\leqslant x-1\)，从而
\(x-1-\ln x\geqslant0\)，即
\[
\frac{\e^{x-1}}x=\e^{x-1-\ln x}\geqslant1
\]
若 \(\e x\ln x+2>1\)，则 \(f(x)>1\). 若等号成立，则
\(x=\e^{-1}\)，此时 \(x\ne1\)，故上式中的第一个因子严格大于 \(1\)，仍有
\(f(x)>1\). 因此对任意 \(x>0\)，均有 \(f(x)>1\).
\end{enumerate}
\end{solution}

\section{选考题：解答题}

\begin{problem}
如图，四边形 \(ABCD\) 是 \(\odot O\) 的内接四边形，\(AB\) 的延长线与 \(DC\) 的延长线交于点 \(E\)，且 \(CB = CE\).

\begin{enumerate}
\item 证明：\(\angle D = \angle E\)；
\item 设 \(AD\) 不是 \(\odot O\) 的直径，\(AD\) 的中点为 \(M\)，且 \(MB = MC\)，证明：\(\triangle ADE\) 为等边三角形.
\end{enumerate}

\FigureLayoutDeclare{img/2014/new_curriculum_paper_1_science/q22_fig1.png}{6.5cm}{102.712bp 48.476bp 0bp 11.039bp}
\bitmapfigure[float=inline,max width=0.62\linewidth]{img/2014/new_curriculum_paper_1_science/q22_fig1.png}
\FloatBarrier
\end{problem}

\begin{solution}
\begin{enumerate}
\item 因为 \(CB=CE\)，所以三角形 \(BCE\) 为等腰三角形，故
\(\angle CBE=\angle BEC=\angle E\). 又因为 \(A\)，\(B\)，\(E\) 三点共线，
\(\angle CBE\) 与四边形 \(ABCD\) 在 \(B\) 点的内角互为补角. 对圆内接四边形，
\(\angle ABC+\angle ADC=180^{\circ}\)，所以
\[
\angle ADC=\angle CBE=\angle BEC=\angle E
\]
即 \(\angle D=\angle E\).

\item 因为 \(AD\) 不是直径，圆心 \(O\) 与弦 \(AD\) 的中点 \(M\) 不重合，且
\(OM\perp AD\). 又由 \(OB=OC\) 和 \(MB=MC\)，直线 \(OM\) 是线段 \(BC\) 的
垂直平分线，所以 \(OM\perp BC\). 从而
\(AD\parallel BC\).

由 \(EA\) 与 \(EB\) 共线、\(AD\parallel BC\)，得
\(\angle EAD=\angle EBC\). 因为 \(CB=CE\)，三角形 \(BCE\) 为等腰三角形，
所以 \(\angle EBC=\angle BEC=\angle E\). 因而
\(\angle EAD=\angle E\). 第（1）问又给出
\(\angle ADE=\angle D=\angle E\). 于是三角形 \(ADE\) 的三个内角均等于
\(\angle E\)，故
\(3\angle E=180^{\circ}\)，即三个角均为 \(60^{\circ}\)，所以
\(\triangle ADE\) 为等边三角形.
\end{enumerate}
\end{solution}

\begin{problem}
已知曲线 \(C: \frac{x^2}{4} + \frac{y^2}{9} = 1\)，直线 \(l: \begin{cases}
x = 2 + t,\\
y = 2 - 2t,
\end{cases}\) （\(t\) 为参数）.

\begin{enumerate}
\item 写出曲线 \(C\) 的参数方程，直线 \(l\) 的普通方程；
\item 过曲线 \(C\) 上任一点 \(P\) 作与 \(l\) 夹角为 \(30^{\circ}\) 的直线，交 \(l\) 于点 \(A\)，求 \(|PA|\) 的最大值与最小值.
\end{enumerate}
\end{problem}

\begin{solution}
\begin{enumerate}
\item 令 \(x=2\cos\theta\)，\(y=3\sin\theta\)，其中
\(0\leqslant\theta<2\pi\)，则
\[
\left\{
\begin{aligned}
x&=2\cos\theta,\\
y&=3\sin\theta.
\end{aligned}
\right.
\]
是曲线 \(C\) 的参数方程. 由直线的参数方程消去 \(t\)，得
\(t=x-2\)，从而 \(y=2-2(x-2)=6-2x\)，所以直线 \(l\) 的普通方程为
\(2x+y-6=0\).

\item 设点 \(P=(2\cos\theta,3\sin\theta)\). 点 \(P\) 到直线 \(l\) 的距离为
\[
d=\frac{|4\cos\theta+3\sin\theta-6|}{\sqrt5}
\]
因为 \(4\cos\theta+3\sin\theta\leqslant5<6\)，所以
\[
d=\frac{6-4\cos\theta-3\sin\theta}{\sqrt5}
\]
直线 \(PA\) 与 \(l\) 的夹角为 \(30^{\circ}\)，故
\(d=|PA|\sin30^{\circ}=\dfrac12|PA|\)，即 \(|PA|=2d\).

由柯西不等式，
\[
-5\leqslant4\cos\theta+3\sin\theta\leqslant5
\]
且两端均可取到. 因而
\[
\frac{2}{\sqrt5}\leqslant|PA|\leqslant\frac{22}{\sqrt5}
\]
所以 \(|PA|\) 的最大值为 \(\dfrac{22\sqrt5}{5}\)，最小值为
\(\dfrac{2\sqrt5}{5}\).
\end{enumerate}
\end{solution}

\begin{problem}
若 \(a > 0\)，\(b > 0\)，且 \(\frac{1}{a} + \frac{1}{b} = \sqrt{ab}\).

\begin{enumerate}
\item 求 \(a^3 + b^3\) 的最小值；
\item 是否存在 \(a\)，\(b\)，使得 \(2a + 3b = 6\)？并说明理由.
\end{enumerate}
\end{problem}

\begin{solution}
\begin{enumerate}
\item 由题设乘以 \(ab>0\)，得
\[
a+b=(ab)^{3/2}
\]
令 \(t=\sqrt{ab}>0\)，则 \(a+b=t^3\). 又由基本不等式
\(a+b\geqslant2\sqrt{ab}=2t\)，得 \(t\geqslant\sqrt2\). 因而
\[
a^3+b^3=(a+b)^3-3ab(a+b)=t^9-3t^5=t^5(t^4-3)
\]
函数 \(t^5(t^4-3)\) 的导数为
\[
3t^4(3t^4-5)
\]
当 \(t\geqslant\sqrt2\) 时，\(3t^4-5\geqslant7>0\)，所以该函数在
\([\sqrt2,+\infty)\) 上递增. 最小值在 \(t=\sqrt2\) 时取得，为
\[
(\sqrt2)^5\bigl((\sqrt2)^4-3\bigr)=4\sqrt2
\]
等号对应 \(a=b=\sqrt2\)，故最小值为 \(4\sqrt2\).

\item 由 \(a+b=(ab)^{3/2}\) 和 \(a+b\geqslant2\sqrt{ab}\)，仍得
\(ab\geqslant2\). 若存在正数 \(a\)，\(b\) 满足 \(2a+3b=6\)，则
\(a=3-\dfrac32b\)，且 \(0<b<2\). 从而
\[
ab=\frac32b(2-b)\leqslant\frac32<2
\]
与 \(ab\geqslant2\) 矛盾. 所以不存在这样的 \(a\)，\(b\).
\end{enumerate}
\end{solution}
